\ifdefined\gkver\else\def\gkver{student}\fi
\documentclass[\gkver]{gaokaozhenti}
\begin{document}

\chapter{2026年湖北}

\begin{enumerate}
\item
%% number: 1
%% paperName: 2026年湖北高考第 1 题 4 分
%% typeId: 1
%% type: 单选题
%% chapter: 原子和原子核
%% point: 人工核反应
%% method: 无
%% score: 4
%% degree: 950
%% duplicateId: 0
%% body:
$ 2025 $ 年 $ 6 $ 月，我国成功掌握商用堆生产钇 $ 90 $（\ce{^{90}_{39} Y}）玻璃微球的技术，可实现批量化生产。在核反应堆中，\ce{^{90}_{39} Y} 由核反应 \ce{^{a}_{39} Y + ^{1}_{b} X -> ^{90}_{39} Y} 生成、则 \xzanswer{D}


\fourchoices[answer=D]
{$ X $ 是质子}
{$ X $ 是电子}
{$ a =90 $}
{$ b =0 $}

%% answer: D
%% memo:
\memoanswer{
核反应遵循质量数守恒和电荷数守恒，对反应式 \ce{^{a}_{39} Y + ^{1}_{b} X -> ^{90}_{39} Y} 分析如下

由质量数守恒 $ a+ 1=90 $ 解得 $ a =89 $

由电荷数守恒 $ 39 +b =39 $ 解得 $ b =0 $

可知 $ X $是中子。

故选 $ D $。
}

\item
%% number: 2
%% paperName: 2026年湖北高考第 2 题 4 分
%% typeId: 1
%% type: 单选题
%% chapter: 光学
%% point: 光电效应
%% method: 无
%% score: 4
%% degree: 850
%% duplicateId: 0
%% body:
光伏产业已发展成为我国新能源三大名片之一。太阳能电池发电时，材料原子中的电子被光子激发（这也是一种光电效应），激发后的电子向电池负极移动、形成一个持续的电动势。关于太阳能电池，下列说法正确的是 \xzanswer{C} 

\fourchoices[answer=C]
{任何频率的光均可使太阳能电池发电}
{太阳光的能量可以全部转化为电池的电能}
{太阳能电池在多云的白天，仍然可以发电}
{太阳能电池发电的功率与太阳光的照射方向无关}

%% answer: C
%% memo:
\memoanswer{
已知光电效应方程为 $E_{k} = h\nu - W_{0}$，其中逸出功 $W = h\nu_{0}$，$\nu_{0}$ 为材料的极限频率。

A．只有入射光频率 $v \geq \nu_{0}$ 时，才能激发出光电子实现发电，低于极限频率的光无法产生光电效应，故 A 错误；

B．太阳能电池的光电转化效率远小于 100\%，太阳光的能量会有部分转化为内能等其他形式的能量，不可能全部转化为电能，故 B 错误；

C．多云的白天仍有频率满足要求的太阳光透过云层照射到电池板，可发生光电效应，因此仍然可以发电，故 C 正确；

D．设垂直于光传播方向的光强为 $ I_{0} $，太阳光入射方向与电池板法线的夹角为 $ \theta $，则电池板单位面积接收的光功率为 $ I = I_{0} \cos \theta $，$ \theta $不同时接收的光功率不同，发电功率也不同，因此发电功率与照射方向有关，故 D 错误。

故选 C。
}

\item
%% number: 3
%% paperName: 2026年湖北高考第 3 题 4 分
%% typeId: 1
%% type: 单选题
%% chapter: 气体性质
%% point: 理想气体状态方程
%% method: 无
%% score: 4
%% degree: 750
%% duplicateId: 0
%% body:
容器中封闭一定质量的理想气体。缓慢改变气体状态，使其体积变为初态的一半，压强变为初态的 $ 3 $ 倍。与气体的初态相比、末态 \xzanswer{B}


\fourchoices[answer=B]
{气体的温度是初态的$ 3 $倍}
{气体分子的平均动能更大}
{每个气体分子的速率都更大}
{单位时间内与容器壁单位面积碰撞的分子数更少}

%% answer: B
%% memo:
\memoanswer{
根据理想气体状态方程 $\frac{p_{1}V_{1}}{T_{1}} = \frac{p_{2}V_{2}}{T_{2}}$。

代入 $p_{2} = 3p_{1}$，$V_{2} = \frac{1}{2}V_{1}$，可得末态温度 $T_{2} = \frac{p_{2}V_{2}}{p_{1}V_{1}}T_{1} = 1.5T_{1}$。

A．末态温度是初态的 $1.5$ 倍，不是 $3$ 倍，故 A 错误；

B．温度是分子平均动能的唯一宏观标志，末态温度更高，气体分子平均动能更大，故 B 正确；

C．分子平均动能增大是统计规律，仅代表分子整体的速率分布向高速区偏移，不代表每个分子的速率都更大，故 C 错误；

D．气体压强由单位时间内与容器壁单位面积碰撞的分子数、分子平均碰撞冲量共同决定：末态温度升高，分子平均碰撞冲量增大，且压强变为初态的 $3$ 倍，因此单位时间碰撞的分子数更多，故 D 错误。

故选 B。
}
\memoanswer[补充]{
单位时间内与容器壁单位面积碰撞的分子数$\Gamma \propto \frac{P}{\sqrt{T}}$
}

\item
%% number: 4
%% paperName: 2026年湖北高考第 4 题 4 分
%% typeId: 1
%% type: 单选题
%% chapter: 电场
%% point: 带电粒子的直线运动
%% method: 无
%% score: 4
%% degree: 550
%% duplicateId: 0
%% body:
一带正电粒子仅受静电力作用，从静止开始沿直线运动，$ 0 \sim 4 t_{0} $ 时间内，其速度 $ v $ 随时间 $ t $ 的变化图像是如图所示的三条线段。$ t_{0} $ 时刻、$ 1.5 t_{0} $ 时刻、$ 3.5 t_{0} $ 时刻粒子所在位置分别为 $ a $ 点、$ b $ 点、$ c $ 点。下列说法正确的是 \xzanswer{C}


\onepicture{TikZ/01.pdf}

\fourchoices[answer=C]
{$b$ 点的电势比 $a$ 点的高}
{$b$ 点的电场强度比 $a$ 点的大}
{粒子在 $a$、$c$ 两点的电势能相等}
{粒子在 $a$、$c$ 两点受到的电场力大小相等}

%% answer: C

%% memo:
\memoanswer{
A．在 $ 0 \sim 2 t_{0} $ 内，粒子做匀加速运动，加速度 $ a_{1} = \frac{v_{0}}{2 t_{0}} $，电场方向与运动方向相同。沿电场线方向电势降低，则 $ b $ 点电势比 $ a $ 点低，故 A 错误。

B．在 $ 0 \sim 2 t_{0} $ 内加速度恒定，$ a $ 点与 $ b $ 点电场强度大小相等，故 B 错误。

C．在 $ t_{0} $ 时刻，速度 $ v_{a} = \frac{v_{0}}{2} $。在 $ 3.5 t_{0} $ 时刻，处于 $ 3 t_{0} \sim 4 t_{0} $ 的匀减速阶段，加速度 $ a_{3} = -\frac{v_{0}}{t_{0}} $，速度 $ v_{c} = v_{0} + a_{3} \cdot 0.5 t_{0} = \frac{v_{0}}{2} $。则 $ a $、$ c $ 两点动能相等，由能量守恒知电势能相等，故 C 正确。

D．设粒子质量为 $ m $，则 $ a $ 点电场力大小 $ F_{a} = m \frac{v_{0}}{2 t_{0}} $，$ c $ 点电场力大小 $ F_{c} = m \frac{v_{0}}{t_{0}} $，两者不相等，故 D 错误。

故选 $ C $。
}
\item
%% number: 5
%% paperName: 2026年湖北高考第 5 题 4 分
%% typeId: 1
%% type: 单选题
%% chapter: 力物体的平衡
%% point: 三力平衡
%% method: 无
%% score: 4
%% degree: 550
%% duplicateId: 0
%% body:
如图所示，半径为 $ R $ 的光滑圆环固定在竖直平面内，原长为 $ 2 \sqrt 3 R $ 的轻质弹簧一端固定在圆环上与圆心 $ O $ 等高的位置，另一端连接一个套在圆环上质量为 $ m $ 的小球，平衡时弹簧与水平方向的夹角为 $ 30 {^\circ} $。弹簧在弹性限度内，重力加速度大小为 $ g $，弹簧的劲度系数为 \xzanswer{C}
\onepicture[w=4.48cm, h=4.48cm]{\includesvg{figs/02}}

\fourchoices[answer=C]
{$ \frac{{\sqrt 3 mg}}{{6R}} $}
{$ \frac{{mg}}{{2R}} $}
{$ \frac{{\sqrt 3 mg}}{{3R}} $}
{$ \frac{{mg}}{R} $}

%% answer: C

%% memo:
\memoanswer{
根据题意，设此时弹簧的长度为 $ L $，由几何关系有 $ \frac{L}{\sin 120^{\circ}} = \frac{R}{\sin 30^{\circ}} $，解得 $ L = \sqrt{3} R $。

可知此时弹簧被压缩，且形变量为 $ \Delta x = L_{0} - L = \sqrt{3} R $。

对小球受力分析，如图所示。

\onepicture[w=5.5cm]{figs/05_an.png}

由平衡条件有 $ F_{N} \cos 60^{\circ} = k \Delta x \cos 30^{\circ} $，$ F_{N} \sin 60^{\circ} = k \Delta x \sin 30^{\circ} + mg $。

联立解得 $ k = \frac{\sqrt{3}mg}{3R} $。

故选 $ C $。
}
\item
%% number: 6
%% paperName: 2026年湖北高考第 6 题 4 分
%% typeId: 1
%% type: 单选题
%% chapter: 曲线运动
%% point: 开普勒定律
%% method: 无
%% score: 4
%% degree: 450
%% duplicateId: 0
%% body:
已知某卫星绕地球做椭圆运动，在近地点所受的地球引力为其在地面附近的 $ \frac{1}{4} $，在远地点所受的地球引力为其在地面附近的 $ \frac{1}{16} $。地面附近的重力加速度大小为 $ g $，地球半径为 $ R $，该卫星的运动周期为 \xzanswer{B}


\fourchoices[answer=B]
{$ 4 \pi\sqrt {\frac{{2R}}{g}} $}
{$ 6 \pi\sqrt {\frac{{3R}}{g}} $}
{$ 16 \pi\sqrt {\frac{{R}}{g}} $}
{$ 12 \pi\sqrt {\frac{{6R}}{g}} $}

%% answer: B

%% memo:
\memoanswer{
根据题意可知，对于地球表面的物体，由万有引力等于重力有 $ \frac{GMm}{R^{2}} = mg $，可得 $ GM = gR^{2} $。

由万有引力公式有 $ F = \frac{GMm}{r^{2}} $，可知近地点到焦点的距离为 $ 2R $，远地点到焦点的距离为 $ 4R $，则卫星椭圆轨道的半长轴为 $ a = \frac{2R + 4R}{2} = 3R $。

设卫星在轨道半径为 $ R $ 的圆周上做圆周运动的周期为 $ T $，在椭圆轨道上的运动周期为 $ T_{1} $，由开普勒第三定律有 $ \frac{R^{3}}{T^{2}} = \frac{(3R)^{3}}{T_{1}^{2}} $。

由万有引力提供向心力有 $ \frac{GMm}{R^{2}} = m \frac{4\pi^{2}}{T^{2}} R $，联立解得 $ T_{1} = 6\pi\sqrt{\frac{3R}{g}} $。

故选 $ B $。
}
\item
%% number: 7
%% paperName: 2026年湖北高考第 7 题 4 分
%% typeId: 1
%% type: 单选题
%% chapter: 电磁感应
%% point: 电磁感应与动量定理
%% method: 无
%% score: 4
%% degree: 450
%% duplicateId: 0
%% body:
如图所示，固定在水平面内的两足够长平行金属导轨相距 $ L $，其左端用导线连接一阻值为 $ R $ 的电阻，两导轨间存在垂直于导轨平面向下的匀强磁场，磁感应强度大小为 $ B $。质量为 $ m $ 的导体棒跨放在两导轨上，导体棒在两导轨间的电阻为 $ 2 R $。导轨上 $ a $、$ b $、$ c $、$ d $ 是正方形的四个顶点。导体棒以大小为 $ \frac{{4{B^2}{L^3}}}{{mR}} $ 的初速度从 $ ab $ 出发沿导轨向右运动，刚过 $ cd $ 后，$ abcd $ 区域内磁场的磁感应强度方向不变、大小在极短时间内增大到 $ 3 B $，然后保持不变。导体棒运动过程中始终与导轨垂直，且与导轨接触良好，不计导轨电阻与所有摩擦，导体棒最终静止的位置离初始位置的距离为 \xzanswer{B}
\onepicture[w=5.80cm, h=3.48cm]{\includesvg{figs/03}}

\fourchoices[answer=B]
{$ 9 L $}
{$ 10 L $}
{$ 11 L $}
{$ 12 L $}

%% answer: B

%% memo:
\memoanswer{
根据题意，设导体棒最终静止的位置离初始位置的距离为 $ x $，对导体棒从开始运动到静止的过程中，由动量定理有 $ -B \overline{I} L \Delta t = 0 - mv_{0} $。

又有 $ \overline{I} \Delta t = \frac{\overline{E}}{R + 2R} \Delta t = \frac{\Delta\Phi}{3R} $，$ \Delta\Phi = 3BL^{2} + BL(x - L) = 2BL^{2} + BLx $。

联立解得 $ x = 10L $。

故选 $ B $。
}
\item
%% number: 8
%% paperName: 2026年湖北高考第 8 题 4 分
%% typeId: 2
%% type: 多选题
%% chapter: 交流电
%% point: 变压器
%% method: 无
%% score: 4
%% degree: 550
%% duplicateId: 0
%% body:
如图所示，有一理想变压器，其原线圈输入电压为 $ u =100 \sqrt 2 \sin 100 \pi t \UV $，在一个副线圈上接额定电压为 $ 220 \UV $、额定功率为 $ 10 \UW $ 的灯泡甲，在另一个单匝副线圈上接额定电压为 $ 2 \UV $、额定功率为 $ 1 \UW $ 的灯泡乙，两灯泡均能正常发光。下列说法正确的是 \xzanswer{AC}
\onepicture[w=7.38cm, h=3.45cm]{\includesvg{figs/04}}

\fourchoices[answer=AC]
{原线圈的匝数为 $ 50 $}
{输入电压的频率为 $ 100 \UHz $}
{原线圈的输入功率为 $ 11 \UW $}
{流过原线圈的电流最大值为 $ 0.11 \UA $}

%% answer: AC

%% memo:
\memoanswer{
A．原线圈输入电压有效值为 $ U_{1} = \frac{100\sqrt{2}}{\sqrt{2}} \UV = 100 \UV $。设原线圈匝数为 $ n_{1} $，由理想变压器匝数与电压关系有 $ \frac{U_{1}}{U_{3}} = \frac{n_{1}}{1} $，又有 $ U_{3} = 2 \UV $，解得 $ n_{1} = 50 $，故 A 正确；

B．由原线圈输入电压表达式可知 $ \omega = 100\pi \Urads $，则原线圈输入电压的频率为 $ f = \frac{\omega}{2\pi} = 50 \UHz $，故 B 错误；

C．理想变压器原线圈的输入功率等于副线圈输出功率之和为 $ P_{1} = P_{2} + P_{3} = 11 \UW $，故 C 正确；

D．原线圈电流的有效值为 $ I_{1} = \frac{P_{1}}{U_{1}} = 0.11 \UA $，则流过原线圈的电流最大值为 $ I_{m} = \sqrt{2} I_{1} = 0.11\sqrt{2} \UA $，故 D 错误。

故选 $ AC $。
}
\item
%% number: 9
%% paperName: 2026年湖北高考第 9 题 4 分
%% typeId: 2
%% type: 多选题
%% chapter: 曲线运动
%% point: 竖直平面圆周运动
%% method: 无
%% score: 4
%% degree: 350
%% duplicateId: 0
%% body:
如图所示，半径为 $ R $、内壁光滑的圆环轨道固定在竖直平面内。初始时刻，质量为 $ m $ 的小球甲从圆环内表面最高处以大小为 $ v_{0} = \sqrt {2gR} $（$ g $为重力加速度大小）的水平初速度向右运动，同时质量为 $ 2 m $ 的小球乙从圆环内表面最低处以某一水平初速度向左运动。当甲第一次运动到圆环最低点时，乙恰好第一次运动到圆环最高点。不计空气阻力，下列说法正确的是 \xzanswer{BD}
\onepicture[w=4.48cm, h=6.12cm]{\includesvg{figs/05}}

\fourchoices[answer=BD]
{乙的初速度大小为 $ 3 \sqrt {gR} $}
{甲、乙两小球运动的周期相等}
{任意时刻两小球的连线均过圆环圆心}
{任意时刻两小球对圆环作用力的合力均不为零}

%% answer: BD
%% memo:
\memoanswer{
A．根据题意当甲第一次运动到最低点时，乙恰好第一次运动到圆环最高点，根据对称性可知，乙的初速度大小等于甲到达最低点时的速度大小，故乙在最高点时的速度大小等于 $v_{c} = \sqrt{2gR}$，设乙球在最低处时速度为 $v_{02}$，由机械能守恒得

$-4mgR = \frac{1}{2} \cdot 2mv_{c}^{2} - \frac{1}{2} \cdot 2mv_{02}^{2}$

解得 $v_{02} = \sqrt{6gR}$，故 A 错误；

B．根据题意当甲第一次运动到最低点时，乙恰好第一次运动到圆环最高点，根据对称性可知，球从最高点到最低点和从最低点到最高点的时间相等，故甲乙两球运动的周期相等，故 B 正确；

CD．分析可知，甲从最高点运动到最低点，速度大小从 $\sqrt{2gR}$ 增大到 $\sqrt{6gR}$，乙从最低点到最高点，速度大小从 $\sqrt{6gR}$ 减小到 $\sqrt{2gR}$，从初始时刻取一很短时间 $\Delta t$，甲经过的弧长短，乙经过的弧长长，转过的角度不同，故不可能任意时刻两小球的连线均过圆心；故 C 错误；

D．根据前面分析可知除两球分别在最高点和最低点时两球连线经过圆心，其它位置均不满足两小球的连线过圆心，乙球在最高点和最低点时，分别有 $F_{N2} + 2mg = 2m \frac{\left( \sqrt{2gR} \right)^{2}}{R}$，$F_{N1} - 2mg = 2m \frac{\left( \sqrt{6gR} \right)^{2}}{R}$

解得 $F_{N1} = 14mg$，$F_{N2} = 2mg$

甲球在最低点和最高点时，分别有 $F_{N2}' - mg = m \frac{\left( \sqrt{6gR} \right)^{2}}{R}$，$F_{N1}' + mg = m \frac{\left( \sqrt{2gR} \right)^{2}}{R}$

解得 $F_{N2}' = 7mg$，$F_{N1}' = mg$

根据牛顿第三定律可知乙球在最高点时对圆环的作用力大小为 $2mg$，此时甲球在最低点对圆环的作用力大小为 $7mg$；当乙球在最低点时对圆环的作用力大小为 $14mg$，此时甲球在最高点对圆环的作用力大小为 $mg$，故这两个位置两球对圆环的作用力的合力不为零，结合前面分析其它位置均不满足两小球的连线过圆心，即其它位置两球对圆环的作用力不共线，故合力不可能为零，故对任意时刻两小球对圆环作用力的合力均不为零，故 D 正确。

故选 BD。
}
\memoanswer[补充]{
	首先要证明在圆轨道上运动的时间与质量无关。
	
	以最高点为参考点，与圆心连线转过的角度为$\theta$，角度速为$\omega^{2} R=2g-3g\cos\theta+\frac{v_0}{R}$，角速度与质量无关，运动时间也与质量无关。
}

\item
%% number: 10
%% paperName: 2026年湖北高考第 10 题 4 分
%% typeId: 2
%% type: 多选题
%% chapter: 曲线运动
%% point: 斜抛运动
%% method: 极值
%% score: 4
%% degree: 350
%% duplicateId: 0
%% body:
某山沟竖直截面图如图所示，山沟的一侧竖直，另一侧是以 $ O $ 点为圆心、$ R $ 为半径的 $ \frac{1}{4} $ 圆弧，圆弧最高点与 $ O $ 点等高。救援队从 $ O $ 点以大小为 $ v_{0} = \frac{{\sqrt {2gR} }}{2} $ 的初速度向该山沟投掷救援物资，其中 $ g $ 是重力加速度大小。物资可视为质点，不计空气阻力。为避免损坏救援物资，要求物资落到圆弧上的速率最小，则物资 \xzanswer{AD}
\onepicture[w=3.95cm, h=3.77cm]{\includesvg{figs/06}}

\fourchoices[answer=AD]
{在空中运动的时间为 $ \sqrt {\frac{{2R}}{g}} $}
{与水平方向成 $ 45 {^\circ} $ 角斜上抛}
{抛出点与落点的高度差为 $ \frac{{\sqrt 2 R}}{2} $}
{落到圆弧上的最小速率为 $ \frac{{\sqrt {6gR} }}{2} $}

%% answer: AD
%% memo:
\memoanswer{
ACD．设落点与 $O$ 点的竖直高度为 $h$，水平位移为 $x$，初速度与水平方向的夹角为 $\theta$，将初速度沿水平和竖直方向分解，可得 $-h = v_{0} \sin \theta \cdot t - \frac{1}{2}gt^{2}$，$x = v_{0} \cos \theta \cdot t$

同时有 $h^{2} + x^{2} = R^{2}$

联立可得 $h = \frac{1}{4}gt^{2} + \frac{R^{2}}{g} \cdot \frac{1}{t^{2}} - \frac{1}{2}R$

设落到圆弧上的速度为 $v$，根据机械能守恒 $mgh + \frac{1}{2}mv_{0}^{2} = \frac{1}{2}mv^{2}$

解得 $v = \sqrt{v_{0}^{2} + 2gh}$

故可知 $h$ 越小，$v$ 越小，故当 $\frac{1}{4}gt^{2} = \frac{R^{2}}{g} \cdot \frac{1}{t^{2}}$ 时，$h$ 取最小值，落到圆弧上的速度最小；

解得 $t = \sqrt{\frac{2R}{g}}$，$h_{\min} = \frac{1}{2}R$，$v_{\min} = \frac{\sqrt{6gR}}{2}$，故 AD 正确，C 错误；

B．根据前面分析，当 $t = \sqrt{\frac{2R}{g}}$，$h_{\min} = \frac{1}{2}R$ 时，代入解得 $\sin \theta = \frac{1}{2}$，$\theta = 30^{\circ}$，故 B 错误。

故选 AD。
}
\memoanswer[补充]{
	可以用函数法，但计算较为繁琐。当高度、初速度大小一定时，以什么角度抛物体抛得最远，最远距离是多少，这些相关结论，也可用于此题的求解。
}

\item
%% number: 11
%% paperName: 2026年湖北高考第 11 题 8 分
%% typeId: 4
%% type: 实验题
%% chapter: 电路
%% point: 测量电阻率
%% method: 无
%% score: 8
%% degree: 550
%% duplicateId: 0
%% body:
在测量铅笔芯电阻率的实验中，实验器材有：待测圆柱形铅笔芯、电压表 $ \mathrm{V} $、电流表 $ \mathrm{A} $、电阻箱 $ R_{0} $、干电池、刻度尺、螺旋测微器、开关 $ S $、滑片 $ P $ 及导线等。实验电路图如图~\ref{2026湖北07} 所示，$ B $、$ C $ 为铅笔芯的两端，滑片 $ P $ 在 $ B $、$ C $ 间左右滑动。
\twopicture[labelA=2026湖北07, labelB=2026湖北08, wA=6.85cm, hA=5.62cm, wB=3.43cm, hB=5.32cm]
{\includesvg{figs/07}}{\includesvg{figs/08}}
\begin{enumerate}
\item
用螺旋测微器测量铅笔芯的直径 $d$，其示数如图~\ref{2026湖北08} 所示，可知 $d=$ \tkanswer{$0.500$} $\Umm[a]$。
\item
适当选择电阻箱 $ R_{0} $ 的阻值、移动滑片 $ P $ 到某位置、测量 $ B $、$ P $ 间的长度 $ L $，闭合开关 $ S $，记录电压表的示数 $ U $，电流表的示数 $ I $，则该铅笔芯电阻率的表达式为 $ \rho = $ \tkanswer{$\frac{{\pi U{d^2}}}{{4IL}}$}（用 $ U $、$ I $、$ d $ 和 $ L $ 表示）。
\item
多次改变滑片 $ P $ 的位置、记录对应的 $ L $、$ U $、$ I $，由实验数据绘制出的 $ U-IL $ 图如图~\ref{2026湖北09} 所示。由此可得该铅笔芯的电阻率 $ \rho = $ \tkanswer{$6.4\times10^{-6}$} $ \UOm[a] $（$ \pi $ 取 $ 3.14 $，保留两位有效数字）。
\onepicture[num=c, showlabel=false, label=2026湖北09, w=11.07cm, h=8.07cm]{\includesvg{figs/09}}
\item
干电池使用一段时间后，其内阻增大，对电阻率的测量结果\tkanswer{无}（填“有”或“无”）影响。	
\end{enumerate}


%% answer:
\jdanswer{
\begin{enumerate}
	\item 
	$0.500$
	
	\item
	$\frac{{\pi U{d^2}}}{{4IL}}$
	
	\item
	$6.4\times10^{-6}$
	
	\item
	无
\end{enumerate}
}

%% memo:
\memoanswer{
（1）读数为 $ d = 0.5 \Umm + 0.0 \times 0.01 \Umm = 0.500 \Umm $。

（2）根据欧姆定律结合电阻定律可知 $ \frac{U}{I} = \rho \frac{L}{S} $，$ S = \pi \cdot \left( \frac{d}{2} \right)^{2} $，联立解得 $ \rho = \frac{\pi U d^{2}}{4IL} $。

（3）根据 $ \rho = \frac{\pi U d^{2}}{4IL} $ 可知 $ U = \frac{4\rho IL}{\pi d^{2}} $，故 $ U-IL $ 图像的斜率为 $ k = \frac{4\rho}{\pi d^{2}} $，得 $ \frac{4\rho}{\pi d^{2}} = \frac{2.3}{70 \times 10^{-3}} $，代入数值解得 $ \rho = 6.4 \times 10^{-6} \UOm $。

（4）干电池的内阻增大，不考虑电表内阻的影响下，实验中测量的 $ U $ 和通过的电流 $ I $ 还是真实值，对实验测量结果无影响。
}
\item
%% number: 12
%% paperName: 2026年湖北高考第 12 题 9 分
%% typeId: 4
%% type: 实验题
%% chapter: 运动和力
%% point: 验证牛顿第二定律
%% method: 无
%% score: 9
%% degree: 450
%% duplicateId: 0
%% body:
某同学设计了一个验证牛顿第二定律的实验，使用的器材有：直导轨、小车、光电门、数字毫秒计、遮光条、天平、钩码、刻度尺、游标卡尺、定滑轮、细绳等。实验装置如图所示，光电门 $ 1 $、$ 2 $ 分别固定在导轨 $ A $、$ B $ 处，遮光条安装在小车上，连接小车两端的细绳分别跨过左右定滑轮、各悬挂 $ 10 $ 个质量均为 $ m_{0} $ 的钩码，定滑轮间的细绳与导轨平行。重力加速度大小为 $ g $。实验步骤如下：
\onepicture[w=13.18cm, h=5.93cm]{\includesvg{figs/10}}
\begin{enumerate}[label=\roman*.]
	\item 
	用天平测量小车和遮光条的总质量 $ m $，用刻度尺测量 $ A $、$ B $ 间的距离 $ s $，用游标卡尺测量遮光条的宽度 $ d $。
	\item 
	调节导轨倾斜程度，使小车沿导轨向右下滑时通过两个光电门的遮光时间近似相等。
	\item 
	从左侧细绳取下 $ n $ 个钩码挂到右侧细绳，再从导轨左端静止释放小车，用数字毫秒计测量遮光条经过光电门 $ 1 $、$ 2 $ 的遮光时间 $ t_{1} $、$ t_{2} $。
\end{enumerate}
回答下列问题：
\begin{enumerate}
	\item 
	步骤 \ronum{2} 操作的主要目的是\tkanswer{平衡摩擦力}。
	
	\item
	由牛顿第二定律，小车的加速度大小 $ a =$\tkanswer{$\frac{{2n{m_0}g}}{{m + 20{m_0}}}$}（用 $ m_{0} $、$ m $、$ g $ 和 $ n $ 表示）。
	
	\item
	步骤 \ronum{3} 中测得的小车加速度大小 $ a =$\tkanswer{$\frac{{{d^2}}}{{2s}}\left( {\frac{1}{{t_2^2}} - \frac{1}{{t_1^2}}} \right)$}（用 $ d $、$ s $、$ t_{1} $ 和 $ t_{2} $ 表示）。
	
	\item
	该同学根据测出的多组 $ a $、$ n $ 值，作出 $ a-n $ 图像。若该图像近似是一条通过原点的直线，则可以验证：质量一定时，物体的加速度与所受合外力成\tkanswer{正比}。
	
	\item
	若 $ a-n $ 图像明显不过原点，原因可能是\tkanswer{B}（单选，填标号）。
	\threechoices[answer=B]
	{步骤 \ronum{3} 中小车没有从静止释放}
	{步骤 \ronum{2} 中导轨的倾斜程度调节不恰当}
	{钩码质量没有远小于小车和遮光条的总质量}
	
\end{enumerate}


%% answer:
\jdanswer{
\begin{enumerate}
	\item 
	平衡摩擦力
	
	\item
	$\frac{{2n{m_0}g}}{{m + 20{m_0}}}$
	
	\item
	$\frac{{{d^2}}}{{2s}}\left( {\frac{1}{{t_2^2}} - \frac{1}{{t_1^2}}} \right)$
	
	\item
	正比
	
	\item
	B
\end{enumerate}
}

%% memo:
\memoanswer{
（1）通过两个光电门的遮光时间近似相等，可知速度相等，小车做匀速运动，故步骤 ii 的主要目的是平衡摩擦力；

（2）根据牛顿第二定律，对整体 $ [10 + n - (10 - n)] m_{0}g = (m + 20m_{0})a $，解得小车的加速度大小为 $ a = \frac{2nm_{0}g}{m + 20m_{0}} $；

（3）根据运动学公式有 $ \left( \frac{d}{t_{2}} \right)^{2} - \left( \frac{d}{t_{1}} \right)^{2} = 2as $，解得 $ a = \frac{d^{2}}{2s} \left( \frac{1}{t_{2}^{2}} - \frac{1}{t_{1}^{2}} \right) $；

（4）根据前面分析有 $ a = \frac{2m_{0}g}{m + 20m_{0}} n $，可知当 $ a-n $ 图像为一条过原点的直线，当质量一定时，物体的加速度与所受到的合外力成正比；

（5）A．小车没有从静止释放，不影响测量的加速度大小，故不会使 $ a-n $ 图像不过原点，故 A 错误；
B．导轨的倾斜程度调节不恰当时，可知平衡摩擦力不足或者过大，根据牛顿第二定律，对整体 $ [10 + n - (10 - n)] m_{0}g + f = (m + 20m_{0})a $，得 $ a = \frac{2nm_{0}g}{m + 20m_{0}} + \frac{f}{m + 20m_{0}} $，图像不过原点，故 B 正确；
C．根据前面分析，对整体分析得出加速度表达式，当钩码质量没有远小于小车和遮光条的总质量，对本实验没有影响，故 C 错误。

故选 $ B $。
}
\item
%% number: 13
%% paperName: 2026年湖北高考第 13 题 9 分
%% typeId: 6
%% type: 计算题
%% chapter: 机械波
%% point: 波的图象
%% method: 无
%% score: 9
%% degree: 550
%% duplicateId: 0
%% body:
一列简谐横波沿 $ x $ 轴传播，波速为 $ v =2 \Ums $。$ t =0 $ 时刻的波形图如图所示，此时 $ x =3\Um $ 处的质点$ P $的振动方向沿 $ y $ 轴负方向，$ x =5 \Um $ 和 $ x =11 \Um $ 处的质点均处于平衡位置。
\begin{enumerate}
	\item 
	求波的传播方向、波长 $ \lambda $ 和周期 $ T $。
	
	\item
	从 $ t =0 $ 时刻开始，经多长时间质点 $ P $ 第二次到达平衡位置？
\end{enumerate}
\onepicture[w=8.44cm, h=4.40cm, align=right]{\includesvg{figs/11}}
%% answer:
\jdanswer{
\begin{enumerate}
	\item 
	沿 $ x $ 负方向传播，$ \lambda =12 \Um $，$ T =6 \Us $
	
	\item
	$ 4 \Us $
\end{enumerate}
}

%% memo:
\memoanswer{
（1）根据题意 $ t = 0 $ 时刻质点 $ P $ 的振动方向沿 $ y $ 轴负方向，根据同侧法可知波沿 $ x $ 负方向传播，根据图像可知波长 $ \lambda = 2 \times (11 - 5) \Um = 12 \Um $。周期为 $ T = \frac{\lambda}{v} = 6 \Us $。

（2）根据波传播的特性，当 $ t = 0 $ 时 $ x = 5 \Um $ 处质点的振动形式传播到 $ P $ 位置时，质点 $ P $ 第一次到达平衡位置，需要的时间为 $ t_{1} = \frac{\Delta x}{v} = \frac{5 - 3}{2} \Us = 1 \Us $。故质点 $ P $ 第二次到达平衡位置时经过的时间为 $ t_{2} = t_{1} + \frac{1}{2}T = 4 \Us $。
}
\item
%% number: 14
%% paperName: 2026年湖北高考第 14 题 16 分
%% typeId: 6
%% type: 计算题
%% chapter: 磁场
%% point: 带电粒子在组合场中的运动
%% method: 无
%% score: 16
%% degree: 450
%% duplicateId: 0
%% body:
如图所示，在 $ xOy $ 平面内，$ y >0 $ 区域存在匀强电场，电场强度大小为 $ E $、方向沿 $ y $ 轴负方向；在 $ y <0 $ 区域，有一个以 $ O $ 为圆心、$ r $ 为半径的半圆形区域，半圆形区域内既无电场也无磁场，半圆形区域外存在匀强磁场，磁场方向垂直于 $ xOy $ 平面向里。一质量为 $ m $、电荷量为 $ q $ 的带正电粒子从坐标为$( 0 , \frac{{25}}{{12}}r )$的 $ M $ 点静止释放，之后从坐标为$( 2r,0 )$的 $ N $ 点第一次射出磁场。不计重力，求：
\begin{enumerate}
	\item 
	粒子第一次进入磁场时的速度大小；
	
	\item
	磁场的磁感应强度大小；
	
	\item
	粒子第二次射出磁场时的位置坐标。
\end{enumerate}
\onepicture[w=7.12cm, h=7.70cm, align=right]{\includesvg{figs/12}}
%% answer:
\jdanswer{
\begin{enumerate}
	\item 
	$ \sqrt {\frac{{25qEr}}{{6m}}} $
	
	\item
	$ \sqrt {\frac{{8mE}}{{3qr}}} $
	
	\item
	$(0,-r) $
\end{enumerate}
}
%% memo:
\memoanswer{
（1）粒子在匀强电场中做匀加速直线运动，由牛顿第二定律得 $qE = ma$

解得加速度大小为 $a = \frac{Eq}{m}$

粒子从坐标为 $(0, \frac{25}{12}r)$ 的 $M$ 点静止释放，沿 $y$ 轴负方向的位移大小为 $y = \frac{25}{12}r$

由匀变速直线运动速度与位移的关系得 $v_{0}^{2} = 2ay$

联立解得粒子第一次进入磁场的速度大小为 $v_{0} = \sqrt{\frac{25qEr}{6m}}$

（2）粒子第一次进入磁场的坐标为 $(0, -r)$，第一次出磁场的坐标为 $(2r, 0)$，粒子在匀强磁场中的运动轨迹如图所示。

\onepicture[w=5.2cm]{figs/14_an1.png}

设粒子的轨迹半径变为 $r_{0}$，由几何关系得

$r_{0} + \sqrt{r_{0}^{2} - r^{2}} = 2r$

解得 $r_{0} = \frac{5}{4}r$

粒子做匀速圆周运动，由洛伦兹力提供向心力，有

$qv_{0}B = \frac{mv_{0}^{2}}{r_{0}}$

解得 $B = \sqrt{\frac{8mE}{3qr}}$

$\cos \theta = \frac{r}{r_{0}} = \frac{4}{5}$

（3）设粒子第一次出磁场时速度与 $x$ 轴负方向的夹角为 $\theta$，由几何关系得

$\sin \theta = \sqrt{1 - \cos^{2} \theta} = \frac{3}{5}$

则

粒子沿 $x$ 轴和 $y$ 轴的速度分别为 $v_{x} = -v_{0} \cos \theta$，$v_{y} = v_{0} \sin \theta$

解得 $v_{x} = -\frac{4v_{0}}{5}$，$v_{y} = \frac{3}{5}v_{0}$

粒子沿 $y$ 轴方向做类竖直上抛运动，由竖直上抛运动规律，可得运动时间 $t = \frac{2v_{y}}{a}$

沿 $x$ 轴方向做匀速直线运动，有 $x_{0} = v_{x}t$

联立解得沿 $x$ 轴方向的位移为 $x_{0} = -4r$

粒子第二次进入磁场时，$x$ 坐标为 $x = 2r + x_{0} = -2r$

沿 $x$ 轴方向的速度为 $v_{x}' = v_{x}$

由运动的对称性，可知沿 $y$ 轴方向的速度为 $v_{y}' = -v_{y}$

则粒子与 $x$ 轴负方向的夹角仍满足 $\left| \frac{v_{y}}{v_{x}} \right| = \left| \frac{v_{y}'}{v_{x}'} \right| = \tan \theta$

因此粒子两次在磁场中的运动轨迹关于 $y$ 轴对称，如图所示。

\onepicture[w=4.8cm]{figs/14_an2.png}

则粒子第二次出磁场的坐标，即为粒子第一次进入磁场的坐标 $(0, -r)$。
}
\memoanswer[补充]{
	半圆顶点与$N$连线与$x$轴夹角为$\theta$，粒子第一次出磁场时与$x$轴的夹角$\sin2\theta= \frac{4}{5} $，画出第二次进入磁场的轨迹图，其与$y$轴相切，切点刚好是半圆顶点。
}


\item
%% number: 15
%% paperName: 2026年湖北高考第 15 题 18 分
%% typeId: 6
%% type: 计算题
%% chapter: 运动和力
%% point: 动量和能量
%% method: 分情况讨论
%% score: 18
%% degree: 350
%% duplicateId: 0
%% body:
在如图所示的竖直平面内，固定在水平地面上的光滑轨道由两倾角均为 $ 30 {^\circ} $ 的足够长轨道与一水平轨道平滑连接而成，连接点分别为 $ A $、$ B $。质量为 $ m $ 的小物块甲放置在左侧倾斜轨道上高 $ h $ 处、质量为 $ km $ 的小物块乙静止在水平轨道上，乙到 $ A $、$ B $ 两点的距离均为 $ h $。现静止释放甲，所有碰撞均为弹性正碰，重力加速度大小为 $ g $，不计空气阻力。
\begin{enumerate}
	\item 
	求甲第一次到达 $ A $ 点时的速度大小。
	
	\item
	求两物块第一次碰撞过程中，乙所受合外力的冲量大小。
	
	\item
	若两物块在水平轨道上发生第二次碰撞，且第二次碰撞前只有一个物块滑上倾斜轨道，求 $ k $ 满足的关系式（不求具体数值）。
\end{enumerate}
\onepicture[w=15.11cm, h=4.01cm, align=right]{\includesvg{figs/13}}
%% answer:
\jdanswer{
\begin{enumerate}
	\item 
	$ v_{0} = \sqrt {2gh} $
	
	\item
	$ I = \frac{2km}{{1 + k}}\sqrt{2gh} $
	
	\item
	$ 0.81 \leq k \leq 1.2 $ 或 $ k \geq 20.1 $
\end{enumerate}
}
%% memo:
\memoanswer{
（1）甲下滑过程轨道光滑，由动能定理：$mgh = \frac{1}{2}mv_{0}^{2}$

解得：$v_{0} = \sqrt{2gh}$

（2）第一次碰撞为弹性正碰，设碰后甲速度为 $v_{1}$，乙速度为 $v_{2}$，取向右为正方向，由动量守恒和机械能守恒：

$mv_{0} = mv_{1} + kmv_{2}$，$\frac{1}{2}mv_{0}^{2} = \frac{1}{2}mv_{1}^{2} + \frac{1}{2}kmv_{2}^{2}$

解得乙弹性碰撞后速度：

$v_{2} = \frac{2v_{0}}{1 + k}$

对乙由动量定理，合外力冲量等于乙动量变化：$I = kmv_{2} - 0$

代入 $v_{0} = \sqrt{2gh}$，得：$I = \frac{2km\sqrt{2gh}}{1 + k}$

（3）由以上分析可得碰后甲速度

$v_{1} = \frac{(1 - k)v_{0}}{1 + k}$

要求第二次碰撞前只有一个物块滑上倾斜轨道且在水平轨道完成碰撞，考虑临界情况：

①碰后甲乙均向右运动，且乙由斜面回到水平轨道时恰好发生第二次碰撞

则有 $v_{1} = \frac{(1 - k)v_{0}}{1 + k} > 0$

可得 $0 < k < 1$

根据运动关系 $\frac{h}{v_{1}} = \frac{h}{v_{2}} + 2\frac{v_{2}}{a}$

其中斜面上的加速度 $a = g \sin \theta$

可得 $k \approx 0.81$

由此临界情况可判断 $k$ 满足的关系式可以为 $0.81 \leqslant k < 1$

②碰后甲静止，乙向右运动后返回碰撞位置再次发生碰撞，满足题意，此时 $k = 1$

③碰后甲向左运动，乙向右运动且乙速度较大，当甲运动到 $A$ 点时，乙从斜面返回恰好撞到甲

则有 $v_{1} = \frac{(1 - k)v_{0}}{1 + k} < 0$

可得 $k > 1$

根据运动关系 $\frac{h}{-v_{1}} = \frac{3h}{v_{2}} + 2\frac{v_{2}}{a}$

解得 $k \approx 1.2$

由此临界情况可判断 $k$ 满足的关系式可以为 $1 < k \leqslant 1.2$

④碰后甲向左运动，乙向右运动且甲速度较大，当甲运动到 $B$ 点时恰好追上乙发生碰撞

则有 $v_{1} = \frac{(1 - k)v_{0}}{1 + k} < 0$

可得 $k > 1$

根据运动关系 $\frac{3h}{-v_{1}} + 2\frac{-v_{1}}{a} = \frac{h}{v_{2}}$

解得 $k \approx 20.1$

由此临界情况可判断 $k$ 满足的关系式可以为 $k \geqslant 20.1$

综合上述四种情况可得 $k$ 满足的关系式 $0.81 \leqslant k \leqslant 1.2$ 或 $k \geqslant 20.1$。
}
\memoanswer[补充]{
	第三问计算较为繁琐，需要讨论几种情况。
	
	$k=1$显然成立。
	
	甲乙同向时，$k<1$，有$\frac{2}{1-k}-\frac{32}{(1+k)^2}>1$，虽然有解析解，但表达式复杂，数值解为$k>0.8136$。
	
	甲乙反向时，需要讨论甲乙谁跑得快，即$1<k<3$和$k>3$，再立不等式方程即可求解。
}

\end{enumerate}

\end{document}
